\slcol{धृतराष्ट्र उवाच ।\\
\Index{धर्मक्षेत्रे कुरुक्षेत्रे} समवेता  युयुत्सवः।\\
मामकाः पाण्डवाश्चैव किमकुर्वत सञ्ज्य  ॥१.१॥}
\translate{
  dhṛtarāṣṭra uvāca\\
  dharmākṣētrē kurukṣētrē samavētā yuyutsavaḥ |\\
  māmakāḥ pāṇḍavāścaiva kimakurvata sañjaya ‖1.1‖}
\cquote{
Dhṛtarāṣṭra said: In the field of \emph{dharma}, at Kurukshetra, gathered and eager for battle, what did mine (my sons or my forces) and the Pandavas do, O Sanjaya?}
\slcol{सञ्जय उवाच ।\\
\Index{दृष्टवा  तु   पाण्डवानीकं}   व्यूढं  दुर्योधनस्तदा ।\\
आचार्यमुपसङ्गम्य  राजा वचनमब्रवीत्  ॥१.२॥}
\translate{sañjaya uvāca\\
dṛṣṭvā tu pāṇḍavānīkaṃ vyūḍhaṃ duryōdhanastadā |\\
ācāryamupasaṅgamya rājā vacanamabravīt ‖1.2‖
}
\cquote{Sanjaya said: Then, seeing the Pandava army arranged in battle formation, King Duryodhana approached his teacher and spoke these words.}


\newpage
\begin{mananam}{\mananamfont Mananam sloka - 1}
\normalfont\mananamtext In my own daily routine of life, when my body-identified tendencies such as desire, anger, fear, jealousy, etc, were challenged by my deeper urge for freedom and the wisdom gathered from the Scriptures and teachers, which force did I choose to follow? What are my battles in daily life where my bad habits and harmful patterns overrule my noble aspirations and resolutions?
\end{mananam}
\WritingHand\enspace{\selfReflect\small{Self Reflection}}
\begin{inspiration}{\mananamfont Inspiration}
\normalfont\mananamtext Be true to yourself and you shall change for the better. Impartial observation is an essential life-skill. It is not enough to have a mere wish to change oneself. One must introspect each day on one’s thoughts, words, and actions, and align with the highest inspiration from sages and saints.
\end{inspiration}
\newpage

\slcol{\Index{पश्यैतां पाण्डुपुत्राणामाचार्य} महतीं चमूम् ।\\
 व्यूढां द्रुपदपुत्रेण तव शिष्येण धीमता ॥१.३॥}
\translate{paśyaitāṃ pāṇḍuputrāṇāmācārya mahatīṃ camūm |\\
vyūḍhāṃ drupadaputrēṇa tava śiṣyēṇa dhīmatā ‖1.3‖
}
\cquote{Behold, O teacher, this mighty army of the sons of Pandu, arrayed in battle formation by the son of Drupada, your wise disciple.}

\slcol{\Index{अत्र शूरा महेष्वासा} भीमार्जुनसमा युधि ।\\
युयुधानो विराटश्च द्रुपदश्च महारथः  ॥१.४॥}
\translate{
  atra  śūrā maheṣvāsā bhīmārjunasamā yudhi |\\    
  yuyudhāno virāṭaśca drupadaśca mahārathaḥ ‖1.4‖  
}
\cquote{Here, heroes wielding mighty bows, equal to Bhīma and Arjuna in battle: Yuyudhāna, Virāṭa, and Drupada, all great chariot-warriors.}
\slcol{\Index{धृष्टकेतुश्चेकितानः} काशिराजश्च वीर्यवान्।\\
पुरुजित् कुन्तिभोजश्च् शैब्यश्च नरपुङ्गवः॥१.५॥}
\translate{
dhṛṣṭaketuścekitānaḥ  kāśirājaśca vīryavān |\\
purujit  kuntibhojaśca śaibyśca narapuṅgavaḥ ‖1.5‖
}
\cquote{Dṛṣṭaketu, Cekitāna, and the valiant King of Kāśī; Purujit, Kuntibhoja, and Śaibya, the foremost among men.}
\slcol{\Index{युधामन्युश्च विक्रान्त} उत्तमौजाश्च वीर्यवान ।\\
सौभद्रो द्रौपदेयाश्च सर्व एव महारथाः ॥१.६॥}
\translate{
yudhāmanyuśca vikrānta uttamaujāśca vīryavān |\\
saubhadrō draupadēyāśca sarva ēva mahārathāḥ ‖1.6‖
}
\cquote{Yudhāmanyu, the valiant; Uttamaujā, the powerful; the son of Subhadrā, and the sons of Draupadī, all are great chariot-warriors.}
\slcol{\Index{अस्माकं तु विशिष्टा} ये तान्निबोध द्विजोत्तं ।\\
नायका मम सैन्यस्य संज्ञार्थं ब्रवीमि ते ॥१.७॥}
\translate{
asmākaṃ tu viśiṣṭā yē tānnibōdha dvijōttama |\\
nāyakā mama sainyasya saṃjñārthaṃ tān bravīmi tē ‖1.7‖
}
\cquote{But know, O best of the twice-born (Dronacharya), the distinguished among us, the leaders of my army; I shall name them for your understanding.}
\slcol{\Index{भवान्भीष्मश्च कर्णश्च} कृपश्च समितिञ्जयः।\\
 अश्वत्थामा विकर्णश्च सौमदत्तिर्जयद्रथः॥१.८॥}
 \translate{
bhavānbhīṣmaśca karṇaśca kṛpaśca samitiñjayaḥ |\\
aśvatthāmā vikarṇaśca saumadattirjayadrathaḥ ‖1.8‖ 
 }
\cquote{You, and also Bhīṣma, Karṇa, and Kṛpa — victorious in battle; Aśvatthāmā, Vikarṇa, and the son of Somadatta (Saumadatti) likewise.}
\slcol{\Index{अन्ये च बहवः} शूरा मदर्थे त्यक्तजीविताः ।\\
 नानाशस्त्रप्रहरणाः सर्वे युद्धविशारदाः ॥१.९॥}
\translate{
anyē ca bahavaḥ śūrā madarthē tyaktajīvitāḥ |\\
nānāśastrapraharaṇāḥ sarvē yuddhaviśāradāḥ ‖1.9‖ 
}
\cquote{And many other heroes too, ready to lay down their lives for my sake, armed with various weapons and missiles, all skilled in the art of war.}
\slcol{\Index{अपर्याप्तं तदस्माकं} बलं भीष्माभिरक्षितम् ।\\
 पर्याप्तं त्विदमेतेषां बलं भीमाभिरक्षितम् ॥१.१०॥}
\translate{
aparyāptaṁ tadasmākaṁ balaṁ bhīṣmābhirakṣitam |\\
paryāptaṁ tvidamētēṣāṁ balaṁ bhīmābhirakṣitam ‖1.10‖ 
 }
\cquote{That force of ours, protected by Bhīṣma, is immeasurable; while this force of theirs, protected by Bhīma, is limited.}
\slcol{\Index{अयनेषु च सर्वेषु} यथाभागमवस्थिताः।\\
 भीष्ममेवाभिरक्षन्तु भवन्तः सर्व एव हि॥१.११॥} 
\translate{
ayanēṣu ca sarvēṣu yathābhāgamavasthitāḥ |\\
bhīṣmamēvābhirakṣantu bhavantaḥ sarva ēva hi ‖1.11‖ 
}
\cquote{Positioned in all strategic posts according to your assigned roles, let every one of you protect Bhīṣma alone.}
\slcol{\Index{तस्य सञ्जनयन्हर्षं} कुरुवृद्धः पितामहः।\\
 सिंहनादं विनद्योच्चैः शङ्खं दध्मौ प्रतापवान्॥१.१२॥}
\translate{
tasya sañjanayanharṣaṁ \\
\hspace*{0.5cm}kuruvṛddhaḥ pitāmahaḥ |\\
siṁhanādaṁ vinadyōccaiḥ śaṅkhaṁ \\
\hspace*{0.5cm}dadhmau pratāpavān ‖1.12‖}
\cquote{(Then) the valiant grandsire (Bhīṣma), eldest of the Kurus, sounding a lion’s roar with great force, blew his conch to rouse Duryodhana’s spirits.}
\slcol{\Index{ततः शङ्खाश्च} भेर्यश्च पणवानकगोमुखाः।\\
 सहसैवाभ्यहन्यन्त स शब्दस्तुमुलोऽभवत्॥१.१३॥}
\translate{
tataḥ śaṅkhāśca bhēryaśca paṇavānakagōmukhāḥ |\\
sahasaivābhyahanyanta sa śabdastumulō:'bhavat ‖1.13‖ 
}
\cquote{Then, conchs and kettle-drums, tabors, trumpets, and cow-horns were all sounded at once, and that tumultuous noise filled the air.}
\slcol{\Index{ततः श्वेतैर्हयैर्युक्ते} महति स्यन्दने स्थितौ।\\
 माधवः पाण्डवश्चैव दिव्यौ शङ्खौ प्रदध्मतुः॥१.१४॥}
\translate{
tataḥ śvētairhayairyuktē mahati syandanē sthitau |\\
mādhavaḥ pāṇḍavaścaiva divyau śaṅkhau pradadhmatuḥ ‖1.14‖}
\cquote{Then, standing in their great chariot drawn by white steeds, Mādhava (Krishna) and the son of Pāṇḍu (Arjuna) sounded their divine conchs.}
\slcol{\Index{पाञ्चजन्यं हृषीकेशो} देवदत्तं धनञ्जयः ।\\
पौण्ड्रं दध्मौ महाशङ्खं भीमकर्मा वृकोदरः ॥१.१५॥}
\translate{
pāñcajanyaṁ hṛṣīkēśō \\
\hspace*{0.5cm}dēvadattaṁ dhanañjayaḥ |\\
pauṇḍraṁ dadhmau mahāśaṅkhaṁ \\
\hspace*{0.5cm}bhīmakarmā vṛkōdaraḥ ‖1.15‖}
\cquote{Hṛṣīkeśa (Krishna) blew Pāñcajanya, Dhanañjaya (Arjuna) sounded Devadatta, and Vṛkodara (Bhima), the doer of mighty deeds, blew his great conch Pauṇḍra.}
\slcol{\Index{अनन्तविजयं राजा} कुन्तीपुत्रो युधिष्ठिरः।\\
नकुलः सहदेवश्च सुघोषमणिपुष्पकौ॥१.१६॥}
\translate{
anantavijayaṁ rājā kuntīputro yudhiṣṭhiraḥ |\\
nakulaḥ sahadēvaśca sughoṣamaṇipuṣpakau ‖1.16‖}
\cquote{King Yudhishthira, son of Kunti, blew Ananta-vijaya; Nakula and Sahadeva sounded Sughoṣa and Maṇipuṣpaka.}
\slcol{\Index{काश्यश्च परमेष्वासः} शिखण्डी च महारथः ।\\
धृष्टद्युम्नो विराटश्च सात्यकिश्चापराजितः ॥१.१७॥}
\translate{
kāśyaśca paramēṣvāsaḥ śikhaṇḍī ca mahārathaḥ |\\
dhṛṣṭadyumnō virāṭaśca sātyakiścāparājitaḥ ‖1.17‖
}
\cquote{The mighty archer, the king of Kāśi; Śikhaṇḍī, the great chariot-warrior; Dhṛṣṭadyumna, Virāṭa, and Sātyaki, undefeated in battle.}
\slcol{\Index{द्रुपदो द्रौपदेयाश्च} सर्वशः पृथिवीपते।\\
सौभद्रश्च महाबाहुः शङ्खान्दध्मुः पृथक्पृथक्॥१.१८॥}
\translate{
drupadō draupadēyāśca\\
\hspace*{0.5cm}sarvaśaḥ pṛthivīpatē |\\
saubhadraśca mahābāhuḥ\\
\hspace*{0.5cm}śaṅkhāndadhmuḥ pṛthakpṛthak ‖1.18‖}
\cquote{Drupada and the sons of Draupadī, all of them, and the mighty-armed son of Subhadrā, each blew their conch separately, O Lord of the Earth (Dhṛtarāṣṭra).}
\slcol{\Index{स घोषो धार्तराष्ट्राणां} हृदयानि व्यदारयत् ।\\
नभश्च पृथिवीं चैव तुमुलो व्यनुनादयन् ॥१.१९॥}
\translate{
sa ghōṣō dhārtarāṣṭrāṇāṁ hṛdayāni vyadārayat |\\
nabhaśca pṛthivīṁ caiva tumulō vyanunādayan ‖1.19‖
}
\cquote{That tumultuous sound (of conches) rent the hearts of Dhṛtarāṣṭra’s sons, and echoed across both sky and earth.}
\slcol{\Index{अथ व्यवस्थितान्दृष्ट्वा} धार्तराष्ट्रान्कपिध्वजः ।\\
प्रवृत्ते शस्त्रसम्पाते धनुरुद्यम्य पाण्डवः ॥१.२०॥\\
\Index{हृषीकेशं तदा} वाक्यमिदमाह महीपते ।}
\translate{
atha vyavasthitān dṛṣṭvā dhārtarāṣṭrān kapidhvajaḥ |\\
pravṛttē śastrasampātē dhanurudyamya pāṇḍavaḥ ‖1.20‖\\
hṛṣīkēśaṁ tadā vākyamidamāha mahīpatē|
}
\cquote{Then, seeing the sons of Dhṛtarāṣṭra arrayed in position, as the clash of weapons was about to begin, the son of Pāṇḍu (Arjuna), whose banner bore Hanumān, raised his bow and spoke these words to Hṛṣīkeśa (Krishna), O Lord of the Earth (Dhṛtarāṣṭra).}
\slcol{अर्जुन उवाच —\\
 सेनयोरुभयोर्मध्ये रथं स्थापय मेऽच्युत ॥ १.२१ ॥\\
\Index{यावदेतान्निरीक्षेऽहं} योद्धुकामानवस्थितान् ।\\
 कैर्मया सह योद्धव्यमस्मिन्रणसमुद्यमे ॥ १.२२ ॥}
\translate{
arjuna uvāca — \\
sēnayōrubhayōrmadhyē rathaṁ sthāpaya mē:'cyuta ‖1.21‖\\
yāvadētānnirīkṣē:'haṁ yōddhukāmānavasthitān |\\
kairmayā saha yōddhavyamasminraṇasamudyamē ‖1.22‖
}
\cquote{Arjuna said: Place my chariot, O Achyuta (Krishna), in between the two armies, so that I may observe those assembled for battle on this eve of war those with whom I must engage with.}
\slcol{\Index{योत्स्यमानानवेक्षेऽहं} य एतेऽत्र समागताः ।\\
धार्तराष्ट्रस्य दुर्बुद्धेर्युद्धे प्रियचिकीर्षवः ॥१.२३॥}
\translate{
yōtsyamānānavēkṣē:'haṁ ya ētē:'tra samāgatāḥ |\\
dhārtarāṣṭrasya durbuddhēryuddhē priyacikīrṣavaḥ ‖1.23‖
}
\cquote{Let me look upon those who stand here, ready to fight on this battlefield, those who seek to please the evil-minded sons of Dhṛtarāṣṭra.}
\slcol{सञ्जय उवाच —\\
\Index{एवमुक्तो हृषीकेशो} गुडाकेशेन भारत ।\\
सेनयोरुभयोर्मध्ये स्थापयित्वा रथोत्तमम् ॥१.२४॥}
\translate{
sañjaya uvāca —
ēvamuktō hṛṣīkēśō guḍākēśēna bhārata |\\
sēnayōrubhayōrmadhyē sthāpayitvā rathōttamam ‖1.24‖
}
\cquote{Sañjaya said: Thus addressed by Guḍākeśa, O Bhārata (Dhṛtarāṣṭra), Hṛṣīkeśa placed that supreme chariot between the two armies.}
\slcol{\Index{भीष्मद्रोणप्रमुखतः} सर्वेषां च महीक्षिताम् ।\\
उवाच पार्थ पश्यैतान्समवेतान्कुरूनिति ॥१.२५॥}
\translate{
bhīṣmadrōṇapramukhataḥ sarvēṣāṁ ca mahīkṣitām |\\
uvāca pārtha paśyaitānsamavētānkurūniti ‖1.25‖
}
\cquote{Before Bhīṣma, Droṇa, and all the kings assembled, He (Krishna) said: “O Pārtha (Arjuna), behold these Kurus gathered here.”}
\slcol{\Index{तत्रापश्यत्स्थितान्पार्थः} पितॄनथ पितामहान् ।\\
आचार्यान्मातुलान्भ्रातॄन्पुत्रान्पौत्रान्सखींस्तथा ॥१.२६॥}
\translate{
tatrāpaśyatsthitānpārthaḥ pitṝnatha pitāmahān |\\
ācāryānmātulānbhrātṝnputrānpautrānsakhīṁstathā ‖1.26‖
}
\cquote{There, Pārtha saw standing before him: fathers, grandfathers, teachers, maternal uncles, brothers, sons, grandsons, and friends;}
\slcol{\Index{श्वशुरान्सुहृदश्चैव}सेनयोरुभयोरपि ।\\
 तान्समीक्ष्य स कौन्तेयः सर्वान्बन्धूनवस्थितान् ॥१.२७॥\\
\Index{कृपया परयाविष्टो} विषीदन्निदमब्रवीत् ।}
\translate{
śvaśurānsuhṛdaścaivasēnayōrubhayōrapi |\\
tānsamīkṣya sa kauntēyaḥ sarvānbandhūnavasthitān ‖1.27‖\\
kṛpayā parayāviṣṭō viṣīdannidamabravīt |
}
\cquote{As well as fathers-in-law and well-wishers, on both sides of the battle-line. Seeing all his kinsmen arrayed before him, the son of Kuntī was overwhelmed with deep compassion and, sorrowing, spoke these words:}
\slcol{अर्जुन उवाच —\\ 
दृष्ट्वेमान्स्वजनान्कृष्ण युयुत्सून्समुपस्थितान् ॥१.२८॥}
\translate{
arjuna uvāca —\\
dṛṣṭvēmānsvajanānkṛṣṇa yuyutsūnsamupasthitān ‖1.28‖
}
\cquote{Arjuna said: Seeing my own kinsmen, O Krishna, arrayed here, intent on battle; and my mouth is parched.}
\slcol{\Index{सीदन्ति मम गात्राणि} मुखं च परिशुष्यति ।\\
वेपथुश्च शरीरे मे रोमहर्षश्च जायते ॥१.२९॥}
\translate{
sīdanti mama gātrāṇi mukhaṁ ca pariśuṣyati |\\
vēpathuśca śarīrē mē rōmaharṣaśca jāyatē ‖1.29‖
}
\cquote{My limbs falter, and my mouth is parched; my body trembles, and my hair stands on end.}

\slcol{\Index{गाण्डीवं स्रंसते} हस्तात्त्वक्चैव परिदह्यते ।\\
न च शक्नोम्यवस्थातुं भ्रमतीव च मे मनः ॥१.३०॥}
\translate{
gāṇḍīvaṁ sraṁsatē hastāttvakcaiva paridahyatē |\\
na ca śaknōmyavasthātuṁ bhramatīva ca mē manaḥ ‖1.30‖
}
\cquote{The Gāṇḍīva slips from my hand, and my skin burns all over. I am unable to remain steady, and my mind seems to reel.}

\slcol{\Index{निमित्तानि च पश्यामि} विपरीतानि केशव ।\\
 न च श्रेयोऽनुपश्यामि हत्वा स्वजनमाहवे ॥१.३१॥}
\translate{
nimittāni ca paśyāmi viparītāni kēśava |\\
na ca śrēyō:'nupaśyāmi hatvā svajanamāhavē ‖1.31‖
}
\cquote{And I see adverse omens, O Keśava (Krishna). I see no good in slaying my own kinsmen in this battle.}
\slcol{\Index{न काङ्क्षे विजयं} कृष्ण न च राज्यं सुखानि च ।\\
किं नो राज्येन गोविन्द किं भोगैर्जीवितेन वा ॥१.३२॥}
\translate{
na kāṅkṣē vijayaṁ kṛṣṇa na ca rājyaṁ sukhāni ca |\\
kiṁ nō rājyēna gōvinda kiṁ bhōgairjīvitēna vā ‖1.32‖
}
\cquote{I desire no victory, O Krishna, nor kingdom, nor pleasures. Of what use to us is kingdom, O Govinda, or enjoyments, or even life itself?}

\slcol{\Index{येषामर्थे काङ्क्षितं} नो राज्यं भोगाः सुखानि च।\\
त इमेऽवस्थिता युद्धे प्राणांस्त्यक्त्वा धनानि च॥१.३३॥}
\translate{
yēṣāmarthē kāṅkṣitaṁ nō rājyaṁ bhōgāḥ sukhāni ca |\\
ta imē:'vasthitā yuddhē prāṇāṁstyaktvā dhanāni ca ‖1.33‖
}
\cquote{For whose sake we desired kingdom, enjoyments, and happiness, they now stand here prepared for battle, having renounced life and wealth.}


\newpage
\begin{mananam}{\mananamfont Mananam sloka - 30}
\mananamtext Let me reflect on those moments in life when I had to face tremendous anxiety and was overwhelmed within by my outward circumstances. In such situations, was I aware that my physical condition was paralysed due to my mental fear? How do I deal with anxiety and fear in my life?
\end{mananam}
\WritingHand\enspace{\selfReflect\small{Self Reflection}}
\begin{inspiration}{\mananamfont Inspiration}
\mananamtext Remain vigilant over your thoughts. Your mental states have an impact on your body. Master some simple yogic breathing to help you free the mind during stress in daily life. 
\end{inspiration}
\newpage

\slcol{\Index{आचार्याः पितरः} पुत्रास्तथैव च पितामहाः।\\
मातुलाः श्वशुराः पौत्राः स्यालाः सम्बन्धिनस्तथा॥१.३४॥}
\translate{
ācāryāḥ pitaraḥ putrāstathaiva ca pitāmahāḥ |\\
mātulāḥ śvaśurāḥ pautrāḥ syālāḥ sambandhinastathā ‖1.34‖
}
\cquote{Teachers, fathers, sons, and likewise grandfathers; maternal uncles, fathers-in-law, grandsons, brothers-in-law, and other kinsmen as well — all stand here prepared for battle.}

\slcol{\Index{एतान्न हन्तुमिच्छामि} घ्नतोऽपि मधुसूदन।\\
अपि त्रैलोक्यराज्यस्य हेतोः किं नु महीकृते॥१.३५॥}
\translate{
ētānna hantumicchāmi ghnatō:'pi madhusūdana |\\ 
api trailōkyarājyasya hētōḥ kiṁ nu mahīkṛtē ‖1.35‖
}
\cquote{These I do not wish to kill, O Madhusūdana (Krishna), even though they strike me down. For the sake of sovereignty over the three worlds, what then to speak of mere earthly rule?}
\slcol{\Index{निहत्य धार्तराष्ट्रान्नः} का प्रीतिः स्याज्जनार्दन ।\\
पापमेवाश्रयेदस्मान्हत्वैतानाततायिनः ॥१.३६॥}
\translate{
nihatya dhārtarāṣṭrānnaḥ kā prītiḥ syājjanārdana |\\
pāpamēvāśrayēdasmānhatvaitānātatāyinaḥ ‖1.36‖
}
\cquote{By killing the sons of Dhṛtarāṣṭra, what joy could be ours, O Janārdana (Krishna)? Sin alone would surely seize us, if we slay these aggressors (those with wrong intentions).}
\slcol{\Index{तस्मान्नार्हा वयं हन्तुं} धार्तराष्ट्रान्सबान्धवान्।\\
स्वजनं हि कथं हत्वा सुखिनः स्याम माधव ॥१.३७॥}
\translate{
tasmānnārhā vayaṁ hantuṁ dhārtarāṣṭrānsabāndhavān |\\
svajanaṁ hi kathaṁ hatvā sukhinaḥ syāma mādhava ‖1.37‖
}
\cquote{Therefore, it isn’t appropriate for us to kill the sons of Dhṛtarāṣṭra, along with their kin. For how could we, by killing our own people, ever be happy, O Mādhava (Krishna)?}

\newpage
\begin{mananam}{\mananamfont Mananam sloka - 37}
\mananamtext Were there times when I justified my actions or inactions due to lack of responsibility? Out of a false sense of sympathy, am I afraid of disassociating with those who pull me downward spiritually and influence me negatively? Am I too weak to say “NO” to people and invitations that are not helpful for my spiritual life?
\end{mananam}
\WritingHand\enspace{\selfReflect\small{Self Reflection}}
\begin{inspiration}{\mananamfont Inspiration}
\mananamtext Rise up to the challenge of life. Take responsibility for your own life. Mercilessly, get rid of all negative associations and environmental influences. Do not surrender your mental peace into the hands of others. Let go of the past and do your best in the present. A bright future is in your grasp.
\end{inspiration}
\newpage

\slcol{\Index{यद्यप्येते न पश्यन्ति} लोभोपहतचेतसः ।\\
कुलक्षयकृतं दोषं मित्रद्रोहे च पातकम् ॥१.३८॥}
\translate{
yadyapyētē na paśyanti lōbhōpahatacētasaḥ |\\
kulakṣayakṛtaṁ dōṣaṁ mitradrōhē ca pātakam ‖1.38‖
}
\cquote{Even if these, with their hearts overpowered by greed, do not see the evil in destroying the family, nor the sin in betraying friends;}
\slcol{\Index{कथं न ज्ञेयमस्माभिः} पापादस्मान्निवर्तितुम्।\\
कुलक्षयकृतं दोषं प्रपश्यद्भिर्जनार्दन॥१.३९॥}
\translate{
kathaṁ na jñēyamasmābhiḥ pāpādasmānnivartitum |\\
kulakṣayakṛtaṁ dōṣaṁ prapaśyadbhirjanārdana ‖1.39‖
}
\cquote{How should we not know to turn away from this sin; we who clearly see the evil caused by destroying the family, O Janārdana (Krishna).}
\slcol{\Index{कुलक्षये प्रणश्यन्ति} कुलधर्माः सनातनाः ।\\
धर्मे नष्टे कुलं कृत्स्नमधर्मोऽभिभवत्युत ॥१.४०॥}
\translate{
kulakṣayē praṇaśyanti kuladharmāḥ sanātanāḥ |\\
dharmē naṣṭē kulaṁ kṛtsnamadharmō:'bhibhavatyuta ‖1.40‖
}
\cquote{When the family is destroyed, the eternal family dharmas perish. With \emph{dharma} lost, the entire family is overcome by \emph{adharma}.}
\slcol{\Index{अधर्माभिभवात्कृष्ण} प्रदुष्यन्ति कुलस्त्रियः ।\\
स्त्रीषु दुष्टासु वार्ष्णेय जायते वर्णसङ्करः ॥१.४१॥}
\translate{
adharmābhibhavātkṛṣṇa praduṣyanti kulastriyaḥ |\\
strīṣu duṣṭāsu vārṣṇēya jāyatē varṇasaṅkaraḥ ‖1.41‖
}
\cquote{When \emph{adharma} prevails, O Krishna, the women of the family become corrupted. When women are corrupted, O Vārṣṇeya (descendant of the Vṛṣṇis, Krishna), there arises intermingling of social orders (\emph{varṇa-saṅkara}).}
\slcol{\Index{सङ्करो नरकायैव} कुलघ्नानां कुलस्य च ।\\
पतन्ति पितरो ह्येषां लुप्तपिण्डोदकक्रियाः ॥१.४२॥}
\translate{
saṅkarō narakāyaiva kulaghnānāṁ kulasya ca |\\
patanti pitarō hyēṣāṁ luptapiṇḍōdakakriyāḥ ‖1.42‖
}
\cquote{Such intermingling of social orders leads only to hell, for the destroyers of the family. Their ancestors fall, deprived of offerings of food and water rites.}
\slcol{\Index{दोषैरेतैः कुलघ्नानां} वर्णसङ्करकारकैः ।\\
उत्साद्यन्ते जातिधर्माः कुलधर्माश्च शाश्वताः ॥१.४३॥}
\translate{
dōṣairētaiḥ kulaghnānāṁ varṇasaṅkarakārakaiḥ |\\
utsādyantē jātidharmāḥ kuladharmāśca śāśvatāḥ ‖1.43‖
}
\cquote{By these faults of the family-destroyers, which bring about the intermingling of social orders, the eternal family dharmas and the duties of lineage are overthrown.}
\slcol{\Index{उत्सन्नकुलधर्माणां} मनुष्याणां जनार्दन ।\\
 नरके नियतं वासो भवतीत्यनुशुश्रुम ॥१.४४॥}
\translate{
utsannakuladharmāṇāṁ manuṣyāṇāṁ janārdana |\\
narakē niyataṁ vāsō bhavatītyanuśuśruma ‖1.44‖
}
\cquote{For those whose family dharmas are destroyed, O Janārdana, perpetual residence in hell is assured. This we have heard from tradition.} 
\slcol{\Index{अहो बत महत्पापं} कर्तुं व्यवसिता वयम् ।\\
 यद्राज्यसुखलोभेन हन्तुं स्वजनमुद्यताः ॥१.४५॥}
\translate{
ahō bata mahatpāpaṁ kartuṁ vyavasitā vayam |\\
yadrājyasukhalōbhēna hantuṁ svajanamudyatāḥ ‖1.45‖
}
\cquote{Alas, what great sin we are resolved to commit. Driven by the greed for the pleasure of a kingdom, we are prepared to slay our own kinsmen.}
\slcol{\Index{यदि मामप्रतीकारमशस्त्रं} शस्त्रपाणयः ।\\
धार्तराष्ट्रा रणे हन्युस्तन्मे क्षेमतरं भवेत् ॥१.४६॥}
\translate{
yadi māmapratīkāramaśastraṁ śastrapāṇayaḥ |\\
dhārtarāṣṭrā raṇē hanyustanmē kṣēmataraṁ bhavēt ‖1.46‖
}
\cquote{Having spoken thus in the midst of battle, Arjuna sat down in the chariot. Casting aside his bow and arrows, his mind overwhelmed with sorrow.}
\newpage
\slcol{सञ्जय उवाच —\\
\Index{एवमुक्त्वार्जुनः सं‍ख्ये} रथोपस्थ उपाविशत् ।\\
विसृज्य सशरं चापं शोकसंविग्नमानसः ॥१.४७॥}
\translate{
ēvamuktvārjunaḥ saṁkhyē rathōpastha upāviśat |\\ 
visṛjya saśaraṁ cāpaṁ śōkasaṁvignamānasaḥ ‖1.47‖
}
\cquote{Sanjaya said: Thus, speaking these words, Arjuna, overwhelmed with sorrow and dejection on the battlefield, took a seat in his chariot and set aside his bow and arrow.}

\begin{center}
\devfont ॐ तत्सदिति श्रीमद्भगवद्गीतासूपनिषत्सु  ब्रह्मविद्यायां योगशास्त्रे\\ श्रीकृष्णर्जुनसंवादे अर्जुनविषादयोगो  नाम प्रथमोऽध्यायः॥\\
\chapEndSlokaTranslate{ōṃ tatsaditi śrīmadbhagavadgītās ūpaniṣatsu \\brahmavidyāyāṃ yōgaśāstrē\\ śrīkṛṣṇārjunasaṃvādē arjunaviṣādayogo\\nāma prathamo'dhyāyaḥ ‖}
\end{center}
